\chapter{刻蚀与 Lift-off：把胶上的图形变成器件}

\begin{chaptergoal}
理解 subtractive etch 与 additive lift-off 的基本选择逻辑；知道侧壁、残胶、过刻、欠刻、redeposition 和金属毛刺为什么会成为电磁与可靠性问题。
\end{chaptergoal}

\section{两条常见路线}

一种路线是先沉整层薄膜，再用光刻胶定义保护区并刻蚀掉不需要的材料；另一种是先形成 resist opening，再沉积材料并 lift-off。二者都能得到图形，但对膜连续性、侧壁、最小间隙、材料兼容性和表面暴露历史的要求不同。

\processchain{图形转移偏差 $\rightarrow$ 实际边缘/侧壁/残留 $\rightarrow$ 电流路径与局部场改变 $\rightarrow$ $L,C,Q_i$, shorts/opens, yield}

\section{不要只看“有没有图形”}

低倍显微镜能确认大结构，但不能替代关键 CD、侧壁和 residue 的抽查。尤其是窄 gap、IDC、coupler 与 meander 转角，应提前定义接受标准和抽样位置。
